%!TEX root=../mythesis.tex
% Chapter 7 -- Conclusion

\chapter{Conclusion and Future Work} % Main chapter title
\chaptermark{Conclusion and Future Work}
\label{ch:conclusion}

This thesis asked what reinforcement learning can do for quantitative trading, where it fails when moved from one trading task to another, and how those failures can be mitigated. It answered the question through three studies spanning second-level position control, five-minute leveraged trading, and daily signal generation. This final chapter summarises the contributions, revisits the comparative perspective that connects them, states their scope and limitations, and sketches the work they invite.

%----------------------------------------------------------------------------------------
\section{Summary of Contributions}
\label{sec:concl_summary}

\paragraph{EarnHFT (\cref{ch:earnhft}).} For second-level high-frequency trading, EarnHFT combines a hindsight, sample-path Q-teacher with trend-conditioned policies and an offline minute-level router. The teacher can be constructed because historical replay makes future prices exogenous to the agent; the corresponding tabular analysis is an idealisation and does not establish optimality for function approximation or an impact-aware live market. In the reported single-seed replays, EarnHFT attains the highest total return on all four datasets, and the component ablation supports the usefulness of value supervision within that experiment.

\paragraph{FineFT (\cref{ch:fineft}).} For leveraged futures, FineFT uses an ensemble TD matrix to assign updates selectively, fits a VAE to each inferred dynamic, and routes by reconstruction-based familiarity scores with a conservative fallback. Leverage magnifies account-return and reward variation rather than the uncertainty of the underlying price transition. The VAE scores are heuristic indicators of distributional familiarity, not calibrated competence boundaries. Across the four reported crypto markets, FineFT has the highest average total return among learned methods and the lowest reported average maximum drawdown, although the evaluation does not isolate the router as the sole cause and is not uniformly dominant by market.

\paragraph{AlphaQD (\cref{ch:alphaqd}).} For formulaic alpha mining, AlphaQD separates a candidate-only GFlowNet reward from a MAP-Elites archive keyed on the normalised daily IC-series shape. This removes archive membership from the reward and provides a structured export set. The chapter's formal results are an abstract equal-weight ICIR lemma and a separate finite-class comparison conditional on a realised archive; neither is a finite-sample guarantee for the fitted ridge portfolio. The method is evaluated on A-share and U.S.\ panels, with its strongest portfolio result on A-shares.

%----------------------------------------------------------------------------------------
\section{The Unifying Perspective Revisited}
\label{sec:concl_perspective}

Across tasks that differ in frequency, instrument, and output, the methods all structure heterogeneity, construct differentiated candidates, and aggregate them. The analogy must be drawn with care. EarnHFT groups historical periods by trend; FineFT assigns transitions to ensemble members from their TD errors; AlphaQD partitions the behaviour space of formulas, not the market. Its cell elites are formulas rather than RL specialists, and a downstream ridge regression combines their outputs. The comparison is therefore an organising vocabulary, not evidence that all three recover locally stationary MDPs.

The narrower connective thread concerns target coupling. EarnHFT trains its router only after freezing the low-level policies; FineFT uses a rule over reconstruction scores; AlphaQD keeps archive membership out of the generator's reward. This motivates a conditional working hypothesis: when aggregation would change a candidate's target as other candidates enter or leave the set, separating the aggregation criterion can stabilise that target and clarify the assumptions behind selection. The studies were neither designed nor executed as a controlled comparison of coupling, so they do not establish a causal performance benefit, and jointly trained mixture-of-experts systems remain valid counterexamples to any universal rule.

There is also concrete methodological reuse. FineFT reuses EarnHFT's hindsight supervision based on dynamic programming and its slope-based segmentation, while adapting both to leveraged futures and ensemble training. AlphaQD then relocates diversity and set selection outside the generator's reward. These are related iterations, but conceptual lineage should not be confused with identical implementations.

%----------------------------------------------------------------------------------------
\section{Limitations}
\label{sec:concl_limitations}

Each study rests on assumptions that bound its claims.

\paragraph{EarnHFT.} The teacher and trend labels use future prices and are available only during historical training. Replay makes prices exactly exogenous to the simulated action, but that is only a small-agent approximation to live trading and cannot support a claim about market impact or optimal order execution. The released learner takes a maximum directly over the target network, so despite the DDQN naming it is not the strict Double-DQN target. The reported runs use one seed (12345); correlated training checkpoints do not provide seed variance, confidence intervals, or significance tests.

\paragraph{FineFT.} The reported experiments use one fixed-seed training run per configuration. The validation period supports several stages---specialist filtering, VAE fitting, score calibration, and router selection---so those choices are not independent validation tests. The VAE reconstruction score is a heuristic familiarity measure rather than a calibrated probability of policy failure, and the supplementary Corn/ES results are much lower-frequency than the main cryptocurrency experiments. These limits do not invalidate the reported within-table ranking, but they bound claims about causal mechanisms and transfer.

\paragraph{AlphaQD.} In the current AIFML configuration reviewed for this thesis, the descriptor is warmed up on 2,000 finite candidates and the PCA/edge model is scheduled for periodic refitting every 100 episodes; refitting reprojects and re-bins all stored elites. The fixed-map theory is therefore an idealised variant, not a guarantee for the moving archive. The AIFML pipeline accepts a seed argument but does not apply it to the main random generators, while PFS creates fresh entropy-seeded generators, so the labelled seed runs are not independently reproducible as controlled seeds. The training loop evaluates the test panel repeatedly, the 20-day labels cross nominal split boundaries, and the U.S.\ and Chinese universes are based on contemporary rather than point-in-time membership. These channels prevent a clean interpretation of the reported test results as untouched out-of-sample estimates.

\paragraph{Common scope.} All three studies are evaluated offline. EarnHFT and FineFT replay recorded data without modelling latency, queue position, partial fills, or the agent's market impact; AlphaQD is a daily signal back-test rather than an execution simulation. The resulting tables compare methods under their respective shared experimental protocols, not live trading performance.

\paragraph{Evaluation protocol.} The test periods are short relative to the range of possible market conditions, validation data are reused for several design and routing choices, and some sensitivity grids are evaluated on the test window. FineFT's AVOL values are reported here on an annualised basis, and its MDD comparison uses the complete final baseline set including RAQR. AlphaQD's archive statistics follow the plotted snapshot consistently. The theoretical results throughout should consequently be read as motivation under explicit premises, not certification of the deployed systems.

%----------------------------------------------------------------------------------------
\section{Future Work}
\label{sec:concl_future}

\paragraph{Testing reward separation in trading.} A direct follow-up is to compare learned, heuristic, and structural aggregation under a controlled position-trading experiment. Such a study could test whether keeping set membership out of a specialist's reward improves target stability; the present three studies do not establish that result.

\paragraph{A unified cross-frequency framework.} The three studies share a vocabulary but not one mathematical construction. A useful framework would make the distinctions explicit---market regime, transition subset, or formula behaviour---and test candidate construction and aggregation independently at each frequency.

\paragraph{Richer descriptors and generators.} AlphaQD's behaviour descriptor could be augmented with additional economically meaningful axes, such as factor-style exposure, and its GFlowNet generator could be replaced by a language-model symbolic policy while preserving the archive-independent reward, so the archive organises diversity rather than becoming a moving component of the reward.

\paragraph{Risk and uncertainty as first-class objectives.} Future work should separate transition noise, model uncertainty, distributional familiarity, and realised account risk rather than treating VAE reconstruction loss as calibrated competence. Point-in-time calibration and proper scoring tests would be required before such signals are used as deployment risk limits.

\paragraph{Toward deployment.} Finally, closing the gap to live trading---modelling latency and market impact, and validating the decoupled-aggregation designs under realistic execution---is the necessary step from these methods to systems that operate in real markets.

\bigskip
\noindent Reinforcement learning is a useful language for the sequential trading problems studied here, but historical performance is credible only when data alignment, account mechanics, selection, and implementation are made auditable. The three studies motivate a working hypothesis about separating aggregation from a candidate's reward when set membership would otherwise move that target. It remains a hypothesis to be tested under controlled and reproducible protocols, not a principle proved by this thesis.
